\exercisesection{对偶化模}

\begin{enumerate}
\def\labelenumi{\arabic{enumi})}
\tightlist
\item
  设 A 为局部 Macaulay 环。
\end{enumerate}

\begin{enumerate}
\def\labelenumi{\alph{enumi})}
\item
  假设 A 具有对偶化模 \(\Omega\)；赋予 A-模 \(\mathrm{A} \oplus \Omega\) 一个 A-代数结构，使 \(\Omega\) 成为平方为零的理想。证明 \(\mathrm{A} \oplus \Omega\) 是 Gorenstein 环（利用极大正则序列化为 A 为 Artin 环的情形；注意 \(\mathrm{A} \oplus \Omega\) 的底座与 A-模 \(\Omega\) 的底座等同）。
\item
  由此推出，为使 A 具有对偶化模，必须且只需存在一个 Gorenstein 环 B 以及从 B 到 A 的满同态。
\end{enumerate}

\begin{enumerate}
\def\labelenumi{\arabic{enumi})}
\setcounter{enumi}{1}
\tightlist
\item
  设 A 为正则局部环，I 为 A 的理想，B 为环 A/I；假设 B 是 Macaulay 环。设 (L,p) 为 A-模 B 的极小投射分解；其长度为 \(c = \text{ht}(I)\)。
\end{enumerate}

\begin{enumerate}
\def\labelenumi{\alph{enumi})}
\item
  设 \(\mathrm{L}^*\) 为复形 \(\operatorname{Homgr}_{\mathrm{A}}(\mathrm{L},\mathrm{A})\)；证明有 \(\mathrm{H}^{i}(\mathrm{L}^{*}) = 0\) 对 \(i\neq c\) 成立，且 B-模 \(\mathrm{H}^c (\mathrm{L}^*)\) 是对偶化的。
\item
  证明下列条件等价：
\end{enumerate}

\begin{enumerate}
\def\labelenumi{(\roman{enumi})}
\item
  A 是 Gorenstein 环；
\item
  复形 L 与 \(L^{*}(-d)\) 同构；
\item
  \(L_{c}\) 是秩为 1 的自由模。
\end{enumerate}

\begin{enumerate}
\def\labelenumi{\arabic{enumi})}
\setcounter{enumi}{2}
\tightlist
\item
  设 A 为具有对偶化模 \(\Omega\) 的环。
\end{enumerate}

\begin{enumerate}
\def\labelenumi{\alph{enumi})}
\item
  为使 \(\Omega\) 同构于 A 的一个理想 I，必须且只需 \(A_{p}\) 是 Gorenstein 环对 A 的每个极小素理想 \(p\) 成立（注意 \(\Omega\) 同构于 \(\bigoplus \Omega_{p}\) 的一个子模）。
\item
  当这些条件满足时，证明 I 等于 A 或者高度为 1，且 A/I 是 Gorenstein 环（先计算 A/I 的深度，再计算 A/I-模 \(\mathrm{Ext}_{\mathrm{A}}^{1}(\mathrm{A / I,I})\)）。
\item
  证明具有对偶化模的因子分解整环是 Gorenstein 环。
\end{enumerate}

\begin{enumerate}
\def\labelenumi{\arabic{enumi})}
\setcounter{enumi}{3}
\item
  设 A 为维数 \(d\) 的局部 Macaulay 环；设 \(\Omega\) 为一个有限生成的 A-模，其内射维数有限、深度为 \(d\)，且使典范同态 \(\mathrm{A} \to \mathrm{End}_{\mathrm{A}}(\Omega)\) 是双射。证明 \(\Lambda\)-模 \(\Omega\) 是对偶化的（对 \(d\) 作归纳，并由 § 3 习题 7 推出：对每个 \(\Omega\)-正则的元素 \(x\)（属于 \(\mathfrak{m}_{\mathrm{A}}\)），典范同态 \(\mathrm{A}/x\mathrm{A} \to \mathrm{End}_{\mathrm{A}/x\mathrm{A}}(\Omega/x\Omega)\) 是双射）。
\item
  设 A 为具有对偶化模 \(\Omega\) 的 Noether 局部环。设 T 为内射维数有限的有限生成 A-模。
\end{enumerate}

\begin{enumerate}
\def\labelenumi{\alph{enumi})}
\item
  证明 \(\mathrm{Ext}_{\mathrm{A}}^{i}(\Omega, \mathrm{T})\) 为零对 \(i > 0\) 成立（利用 § 3 习题 7）。
\item
  证明典范同态 \(\Omega\otimes_{\Lambda}\mathrm{Hom}_{\mathbb{A}}(\Omega,\mathrm{T})\to\mathrm{T}\) 是双射（应用 a) 与 § 8 习题 10）。
\item
  试利用 § 8 习题 8, b) 证明等式 dp \(_{A}\) Hom \(_{A}\) (Ω,T)) = dim(A) - prof(T) 与 prof(Hom \(_{A}\) (Ω,T)) = prof(T)。由此推出，若 prof(T) = dim(A)，则 A-模 T 是若干同构于 Ω 的模的直和。
\end{enumerate}

\(d)\) 设 \(c = \dim(A) - \operatorname{prof}(T)\)。证明存在整数 \(n_0, \ldots, n_c\) 以及一个正合序列

\[
0 \rightarrow \Omega^ {n _ {c}} \rightarrow \dots \rightarrow \Omega^ {n _ {0}} \rightarrow \mathrm{T} \rightarrow 0
\]

（对 c 作归纳，借助 b) 构造一个满同态 \(\Omega^{n_{0}} \to T\)）。

¶ 6) 设 A 为具有对偶化模 Ω 的局部 Noether 环；设 M 为投射维数有限的有限生成 A-模。

\begin{enumerate}
\def\labelenumi{\alph{enumi})}
\item
  证明有 \(\operatorname{Tor}_{i}^{\mathrm{A}}(\Omega, \mathrm{M}) = 0\) 对 \(i > 0\) 成立（对 \(\dim(\mathrm{A})\) 作归纳，考察一个正合序列 \(0 \to \mathrm{N} \to \mathrm{L} \to \mathrm{M} \to 0\)，其中 \(\mathrm{L}\) 是有限型自由模。若 \(x\) 是 \(\mathfrak{m}_{\mathrm{A}}\) 中一个非零因子的元素，则由归纳假设推出 \(\operatorname{Tor}_{i}^{\mathrm{A}}(\Omega, \mathrm{N})\) 当 \(i > 0\) 时为零，由此得 \(\operatorname{Tor}_{i}^{\mathrm{A}}(\Omega, \mathrm{M}) = 0\) 对 \(i > 1\) 成立；又比率为 \(x\) 的乘法映射在 \(\Omega \otimes_{\mathrm{A}} \mathrm{N}\) 中是单射，由此推出每个素理想 \(\mathfrak{p}\)（与 \(\operatorname{Tor}_{1}^{\mathrm{A}}(\Omega, \mathrm{M})\) 相伴）都是 \(\operatorname{Spec}(\mathrm{A})\) 中的极小理想；注意在这样的 \(\mathfrak{p}\) 处，\(\mathrm{A}_{\mathfrak{p}}\)-模 \(\mathrm{M}_{\mathfrak{p}}\) 是自由的）。
\item
  证明 A-模 \(\Omega\otimes_{\mathrm{A}}\mathrm{M}\) 的内射维数有限。
\item
  证明映射 \(\theta: \mathrm{M} \to \mathrm{Hom}_{\Lambda}(\Omega, \Omega \otimes_{\Lambda} \mathrm{M})\)（它由 \(\theta(m)(w) = w \otimes m\) 对 \(m \in \mathrm{M}\)、\(w \in \Omega\) 定义）是同构（借助一个投射分解化为 M 为自由模的情形）。
\end{enumerate}

\(d\) ) 映射 \([\mathrm{T}] \mapsto [\mathrm{Hom}_{\mathrm{A}}(\Omega, \mathrm{T})]\) 与 \([\mathrm{M}] \mapsto [\Omega \otimes_{\mathrm{A}} \mathrm{M}]\) 在内射维数有限的有限生成 A-模的同构类所成的集合与投射维数有限的有限生成 A-模的同构类所成的集合之间定义了互逆的双射，*并且在相应的范畴之间定义了拟逆的范畴等价*。

¶ 7) 设 A 为局部 Macaulay 环，M 为有限生成的 Macaulay A-模。用 \(r(M)\)（相应地，\(t(M)\)）表示 \(\kappa_{A}\)-向量空间 \(\kappa_{A} \otimes_{A} M\)（相应地，\(\operatorname{Ext}_{A}^{p}(\kappa_{A}, M)\)，其中 \(p = \operatorname{prof}(M)\)）的维数，参见 § 1 习题 16。

\begin{enumerate}
\def\labelenumi{\alph{enumi})}
\item
  假设 A 具有对偶化模 \(\Omega\)；设 \(M' = \operatorname{Ext}_A^c(M, \Omega)\)，其中 \(c = \dim(A) - \dim_A(M)\)。证明等式 \(r(M') = t(M)\) 与 \(t(M') = r(M)\)（借助 § 3 命题 7 化为 \(c = 0\) 的情形，再对一个极大割序列取商化为 A 为 Artin 环的情形，并利用 § 8 命题 6）。
\item
  由此推出有 \(t(\mathrm{M}_{\mathfrak{p}}) \leqslant t(\mathrm{M})\) 对每个 \(\mathfrak{p} \in \operatorname{Spec}(\mathrm{A})\) 成立（借助 § 2 习题 2 与 § 1 习题 16 化为 A 完备的情形，并应用 a)）。
\item
  为使模 M 是对偶化的，必须且只需它是忠实的且 \(t(\mathrm{M}) = 1\)（化为 A 完备的情形；应用 a) 与定理 1 的推论）。
\end{enumerate}

\begin{enumerate}
\def\labelenumi{\arabic{enumi})}
\setcounter{enumi}{7}
\tightlist
\item
  设 A 为局部 Gorenstein 环，a 为高度为零的非零理想，b 为 a 的零化子。
\end{enumerate}

\begin{enumerate}
\def\labelenumi{\alph{enumi})}
\item
  证明理想 b 的高度为零，且 A-模 A/b 没有嵌入的伴随素理想（注意与 A/b 相伴的每个素理想都与 A 相伴）。
\item
  假设 A-模 A/α 没有嵌入的伴随素理想；证明 α 是 β 的零化子（化为 dim(A) = 0 的情形，并利用 Matlis 对偶）。
\item
  为使 A/a 是 Macaulay 环，必须且只需 A/b 亦然（利用定理 1 的推论与命题 3 的推论）。当这些条件满足时，A/a-模 b 是对偶化的。
\end{enumerate}

\begin{enumerate}
\def\labelenumi{\arabic{enumi})}
\setcounter{enumi}{8}
\tightlist
\item
  设 \(\Gamma\) 为加法单群 \(\mathbf{N}\) 的一个子单群，使得 \(\mathbf{N} - \Gamma\) 有限；用 \(c\) 表示使 \(\Gamma\) 包含所有整数 \(\geqslant c\) 的最小整数。设 \(k\) 为域，B 为环 \(k[[T]]\)，\(\mathrm{K} = k((\mathrm{T}))\) 为其分式域，A 为 B 的一个子 \(k\)-代数，由元素 \(\mathrm{T}^{\gamma}\)（\(\gamma \in \Gamma\)）生成。
\end{enumerate}

\begin{enumerate}
\def\labelenumi{\alph{enumi})}
\item
  证明 A 是维数 1 的局部整 Noether 环，其分式域为 K，整闭包为 B；理想 c = A : B 等于 BT \(^{c}\)。
\item
  为使 A 是 Gorenstein 环，必须且只需 \(c = 2\operatorname{Card}(\mathbf{N} - \Gamma)\)。
\item
  设 \(\Omega\) 为 K 中由元素 \(T^{-\alpha}\)（\(\alpha \notin \Gamma\)）生成的 A-子模。证明 \(K/\Omega\) 是 Matlis A-模（利用 § 8 命题 2），且 \(\Omega\) 是对偶化 A-模。
\end{enumerate}

\begin{enumerate}
\def\labelenumi{\arabic{enumi})}
\setcounter{enumi}{9}
\tightlist
\item
  设 A 为 Noether 环。设 I 为一个有界内射 A-复形，其同调为有限型的；对每个 A-复形 C，用 D(C) 表示复形 Homgr \(_{A}\) (C,I)。
\end{enumerate}

\begin{enumerate}
\def\labelenumi{\alph{enumi})}
\tightlist
\item
  证明下列条件等价：
\end{enumerate}

\begin{enumerate}
\def\labelenumi{(\roman{enumi})}
\item
  对每个有限生成的 A-模 M，复形的典范态射 \(\boldsymbol{\alpha}_{\mathrm{M}}:\mathrm{M}\to\mathbf{D}(\mathbf{D}(\mathrm{M}))\)（第 5 节）是拟同构；
\item
  对每个同调为有限型的有界 A-复形 C，态射 \(\alpha_{\mathrm{C}}: \mathrm{C} \to \mathbf{D}(\mathbf{D}(\mathrm{C}))\)（见该处）是拟同构；
\item
  典范态射 \(\alpha_{\mathrm{A}}:\mathrm{A}\to \mathrm{Homgr}_{\mathrm{A}}(\mathrm{I},\mathrm{I})\) 是拟同构。
\end{enumerate}

（仿照定理 1 的证明。）

若复形 I 有界、内射，H(I) 有限生成，且满足上述等价条件，我们就说复形 I 是对偶化的。

\begin{enumerate}
\def\labelenumi{\alph{enumi})}
\setcounter{enumi}{1}
\item
  若 A 具有对偶化模 \(\Omega\)，则 \(\Omega\) 的任一有限长度的内射分解 (I, δ) 都定义一个对偶化复形（见该处）。
\item
  设 B 为一个作为有限生成 A-模的 A-代数，I 为一个 A-模的对偶化复形；证明 B-模复形 \(\mathrm{Homgr}_{\mathrm{A}}(\mathrm{B},\mathrm{I})\) 是对偶化的。因此，有限维 Gorenstein 环的每个商环都具有对偶化复形。
\item
  设 I 为有界内射 A-复形。为使 I 是对偶化的，必须且只需对 A 的每个极大理想 m，\(A_{m}\)-模复形 \(I_{m}\) 都是对偶化的。
\item
  设 I 为对偶化 A-复形，P 为秩 1 的投射 A-模，n 为整数。复形 \(1 \otimes_{\mathrm{A}} \mathrm{P}(n)\) 是对偶化的。
\item
  设 I, J 为同调有限型的有界内射 A-复形，\(u: I \rightarrow J\) 为拟同构；为使 J 是对偶化的，必须且只需 I 是对偶化的。
\end{enumerate}

¶ 11) 设 A 为 Noether 环。

\begin{enumerate}
\def\labelenumi{\alph{enumi})}
\tightlist
\item
  假设 \(\Lambda\) 是局部的。设 P, Q 为有界平坦 A-复形。假设 \(\mathrm{H}(\mathrm{P}\otimes_{\mathrm{A}}\mathrm{Q})\) 在次数 \(\neq 0\) 处为零，在次数 0 处是秩 1 的自由模。证明存在一个整数 \(p\in \mathbf{Z}\) 以及拟同构 \(\mathrm{A}\to \mathrm{P}(p)\) 与 \(\mathrm{A}\to \mathrm{Q}(-p)\)。
\end{enumerate}

（利用 A, X 第~66 页的引理 1 与第~62 页的命题 1，化为 P 与 Q 在右边为零的情形；利用该处的命题 1 与 II, § 5, 第 4 节的定理 3，然后构造复形 P' 与 Q' 使 P = A ⊕ P' 且 Q = A ⊕ Q'，并证明有 H(P') = H(Q') = 0。）

\begin{enumerate}
\def\labelenumi{\alph{enumi})}
\setcounter{enumi}{1}
\item
  设 I, J 为对偶化 A-复形（习题 10）。若 A 是局部的，证明存在一个整数 n 以及一个从 J 到 I(n) 上的拟同构（设 P = Homgr \(_A\) (I, J)，Q = Homgr \(_A\) (J, I)；利用第 7 节的态射 \(\nu\) 以及相对于对偶化复形 1 的拟同构 \(\alpha_{J}\) : J → Homgr \(_A\) (Q, I)，构造一个同态 P ⊗ \(_A\) Q → A，并应用 a)）。
\item
  在一般情形，证明存在一个整数 n、一个秩 1 的投射 A-模 L，以及一个从 J 到 \(I \otimes_{A} L(n)\) 上的同态（设 \(L = H(\text{Homgr}_{A}(I, J))\)，并应用 b)）。
\end{enumerate}

\begin{enumerate}
\def\labelenumi{\arabic{enumi})}
\setcounter{enumi}{11}
\tightlist
\item
  设 \(k\) 为域，R 为一个 \(k\)-分次代数，型为 \(\mathbf{N}\)，满足 \(\mathrm{R}_0 = k\) 且 R 是维数 \(d\) 的 Macaulay 环。对每个有限生成的分次 R-模 M，用 \(\mathrm{P}_{\mathrm{M}}(\mathrm{T})\) 表示 Poincaré 级数 \(\sum_{i} (\dim_k \mathrm{M}_i) \mathrm{T}^i\)，并令 \(\mathrm{Q}_{\mathrm{M}}(\mathrm{T}) = (1 - \mathrm{T})^{\dim(M)} \mathrm{P}_{\mathrm{M}}(\mathrm{T})\)；它是 \(\mathbf{Z}[\mathrm{T}, \mathrm{T}^{-1}]\) 中的元素（VIII, § 6, 第 3 节, 命题 5）。
\end{enumerate}

\begin{enumerate}
\def\labelenumi{\alph{enumi})}
\item
  设 \(\Omega\) 为对偶化 R-模；证明 \(\Omega\) 具有一个分次 R-模结构，使得 \(Q_{\Omega}(T) = (-1)^{d} Q_{\mathbb{R}}(T^{-1})\)（把 R 写成多项式代数 \(A = k[X_1, \ldots, X_n]\) 的商，其中 \(\deg(X_i) = d_i\)；若 L 是 A-模 R 的一个分次自由分解，注意复形 \(\mathrm{Homgr}_{\mathrm{A}}(\mathrm{L}, \mathrm{A})(-d)\) 定义了 \(\Omega\) 的一个这样的结构，并利用 VIII, § 4, 第 2 节的例 3）。
\item
  假设 R 是 Gorenstein 环；证明存在整数 \(a \in \mathbf{Z}\) 满足 \(\mathrm{Q}_{\mathbb{R}}(\mathrm{T}^{-1}) = (-1)^{d}\mathrm{T}^{a}\mathrm{Q}_{\mathbb{R}}(\mathrm{T})\)。
\item
  假设 R 是整环，且存在整数 \(a \in \mathbf{Z}\) 使得 \(\mathrm{Q}_{\mathrm{R}}(\mathrm{T}^{-1}) = (-1)^{d}\mathrm{T}^{a}\mathrm{Q}_{\mathrm{R}}(\mathrm{T})\)。证明 R 是 Gorenstein 环（赋予 \(\Omega\) 一个使 \(\mathrm{Q}_{\Omega} = \mathrm{Q}_{\mathrm{R}}\) 的分次结构，并证明 \(\Omega_{0}\) 的一个非零元素构成 \(\Omega\) 的基）。
\item
  设 R 为分次 k-代数 \(k[X,Y]/(X^{3},XY,Y^{2})\)；证明有 \(Q_{R}(T^{-1}) = T^{-2}Q_{R}(T')\)，但 R 不是 Gorenstein 环。
\end{enumerate}

